\addtotoc{Хураангуй}

\vspace{10mm}

\begin{center}
\textbf{\large Хураангуй}
\end{center}

\sloppy

Энэхүү төгсөлтийн судалгааны ажлаар өдөр тутмын тэмдэглэл болон төлөвлөгөөгөө хөтлөх боломжтой “AI Notebook” системийг бүтээв. Систем нь Flutter дээр хөгжүүлсэн гар утасны харагдах хэсэг (frontend), Node.js/Express дээр хөгжүүлсэн REST хэвийн сервер талын хэсэг (RESTful backend), MongoDB өгөгдлийн сан, JWT баталгаажуулалт, мөн OpenRouter API-д суурилсан (OpenAI-compatible gateway) хиймэл оюун ухаанаар стикер үүсгэх модуль (AI sticker generation module)-аас бүрдэнэ. Үүн дээр нэмэлтээр мэдэгдлийн үйлчилгээ (notification service), сануулга хуваарьлалт (reminder scheduling), апп доторх мэдэгдлийн төв (in-app notification center)-ийг нэгтгэн хэрэгжүүлэв.

Судалгаанд Дизайн шинжлэх ухааны судалгааны арга зүй (Design Science Research) ашиглаж, шаардлага тодорхойлох, архитектур боловсруулах, хэрэгжүүлэх, турших гэсэн давталтат үе шатаар ажилласан. Харагдах хэсгийн (frontend) түвшинд хэрэглэгчийн бүртгэл/нэвтрэлт, тэмдэглэл нэмэх-засах-устгах, сэтгэл хөдлөлийн төлөв (mood), календарьтай уялдсан харагдац, профайл удирдлага, мэдэгдлийн дэлгэц, хиймэл оюун ухаанаар стикер үүсгэх интерфейс (AI sticker generation interface) зэрэг боломжуудыг хөгжүүлсэн. Мэдэгдлийн хэсэгт \texttt{reminderAt}/\texttt{stickerUrl} талбар дээр суурилсан нэгтгэлээр мэдэгдлийн жагсаалт үүсгэх, уншсан/уншаагүй төлөв хадгалах, сануулгыг хугацаатай хуваарьлах зэрэг урсгалыг хэрэгжүүлсэн. Сервер талын түвшинд (backend level) хэрэглэгч болон тэмдэглэлийн өгөгдлийн загвар, токеноор хамгаалагдсан API төгсгөлийн цэгүүд (endpoints), тэмдэглэлийн CRUD үйлдэл, зураг үүсгэх интеграц зэрэг үндсэн бизнес логикийг хэрэгжүүлсэн.

Системийн туршилтаар нэвтрэлт, тэмдэглэлийн амьдралын мөчлөг, хиймэл оюун ухаанаар стикер үүсгэх урсгал (AI sticker generation flow), мэдэгдэл ба сануулгын урсгал, токен дээр суурилсан хандалтын хяналт зэрэг гол үйлдлүүд амжилттай ажиллаж байгааг баталгаажуулсан. Үр дүнд нь зөвхөн тэмдэглэл хадгалах бус, хэрэглэгчийн оролцоог нэмэгдүүлэх хиймэл оюунтай интерактив аппликейшний загварыг бодитоор хэрэгжүүлж, цаашдын үүлэн байршуулалт (cloud deployment), алсын түлхэх мэдэгдлийн үйлчилгээ (remote push notification), хэрэглээний шинжилгээ (analytics) зэрэг өргөтгөлийн үндэс тавигдсан.

\vspace{3mm}
\textit{Түлхүүр үгс:} AI Notebook систем, AI sticker generation, Notification service, Calendar integration
